\chapter{Literature Review}
\label{ch:literature}

\section{Design science as the study frame}

Design-science research treats the construction and evaluation of an information-technology artefact as a means of producing knowledge about an identified problem. Hevner et al. distinguish this constructive paradigm from behavioral research and require an artefact, a relevant problem, rigorous evaluation, a research contribution, and effective communication \cite{hevner2004designscience}. Peffers et al. organize design-science research into problem identification, solution objectives, design and development, demonstration, evaluation, and communication \cite{peffers2007dsrm}. These accounts are especially relevant to this thesis because the claimed contribution is not a new foundation model. It is a set of constructs, methods, software components, evidence contracts, and evaluation results assembled around a local research-discovery problem.

The distinction between building and evaluating is important. A route, database table, or source badge demonstrates implementation; it does not establish retrieval effectiveness, correct attribution, or user benefit. Conversely, a high offline metric does not demonstrate that a public interface communicates the metric's limits. This thesis therefore treats the repository and its executed evidence package as the artefact. Demonstration occurs through full-corpus workflows and route probes. Evaluation occurs at several units---corpus records, chunks, questions, claims, citations, profiles, topic labels, review events, and process stages---because a single end-to-end success score would hide where a failure arose.

The design-science frame also disciplines scope. The relevance cycle is the laboratory's need to expose a small, imperfect corpus to students and the public. The rigor cycle draws from scholarly retrieval, RAG evaluation, recommendation, human--AI interaction, and responsible-AI documentation. The design cycle repeatedly connects an observed failure to a design response and a bounded test. For example, stale indexes motivated hash-bound manifests; incomplete answers motivated answer-point review; unverified advice motivated evidence-only and generated modes; and unsafe correction motivated actor-typed, append-only review events. Chapter~\ref{ch:methodology} makes this mapping explicit rather than using ``design science'' as a retrospective label.

\section{Scholarly discovery and research intelligence}

CiteSeer demonstrated autonomous acquisition, citation extraction, and search over online papers \cite{giles1998citeseer}. ArnetMiner combined extraction and mining for expert and topic discovery \cite{tang2008arnetminer}. Semantic Scholar's literature graph, S2ORC, and OpenAlex illustrate complementary graph, full-text-corpus, and open-metadata infrastructures \cite{ammar2018literaturegraph,lo2020s2orc,priem2022openalex}. Their scale enables broad discovery and citation-network analysis, but it also differs from a laboratory collection whose source permissions, missing PDFs, incorrect attachments, and author variants may be inspectable one record at a time.

The TTLAB precursor automated publication collection, display, lay summaries, notifications, and manual podcast preparation, with RAG presented as future work \cite{narine2025automating}. The present artefact retains that public-communication motivation while changing the evidential unit from a publication-list entry or abstract to permitted full-text pages and chunks. It also makes exclusion visible: catalogue representation, local PDF availability, extraction success, eligibility, index coverage, and review state are separate fields. This separation prevents an absent document from disappearing silently and prevents a mismatched PDF from entering evaluation merely because a file exists.

A small system therefore has a different notion of completeness from a large index. It cannot claim global recall, but it can require that every eligible chunk appear in each authoritative representation and that every derived file bind to one ordered corpus snapshot. That form of internal completeness is an engineering claim, not a scholarly-coverage claim. The distinction motivates RQ1 and the manifest-validity checks used throughout the platform.

\section{Information retrieval and evaluation}

BM25 remains a strong lexical baseline \cite{robertson2009bm25}. Feature hashing projects sparse lexical features into a bounded-dimensional space \cite{weinberger2009featurehashing}; because it has no learned language representation, calling it semantic retrieval would confuse implementation convenience with model capability. Sentence-BERT learns sentence embeddings suitable for cosine comparison \cite{reimers2019sentencebert}; Dense Passage Retrieval learns separate query and passage representations for open-domain QA \cite{karpukhin2020dpr}; and SPECTER incorporates citation structure into scholarly-document representations \cite{cohan2020specter}. BEIR shows that zero-shot effectiveness varies materially across datasets and that dense retrieval should not be presumed superior \cite{thakur2021beir}.

These results support three methodological choices. First, lexical, hashed, learned-dense, and hybrid modes are evaluated on identical cases and eligibility boundaries. Second, development tuning is separated from held-out testing. Third, effectiveness is reported with several ranking measures. Set Recall@$k$ answers whether all relevant papers were recovered, reciprocal rank emphasizes the first relevant result, and nDCG incorporates graded relevance and rank discount. None alone measures whether the corpus should have abstained.

Hybrid retrieval is not intrinsically better merely because it combines signals. Recommendation research similarly distinguishes hybridization strategies and their assumptions \cite{burke2002hybrid}. In this platform, the heuristic hybrid combines normalized keyword/vector scores with metadata, section, evidence, topic, and diversity terms; the tuned hybrid changes weights on development cases. Ablation is therefore necessary to determine whether a plausible term helps or harms. Paired tests and confidence intervals are used to resist reading small point-estimate differences as general superiority, while the absence of an equivalence test prevents the inverse claim that two modes are ``the same.''

\section{RAG, attribution, and answer completeness}

RAG conditions generation on retrieved documents \cite{lewis2020rag}. Fusion-in-Decoder aggregates evidence from multiple passages \cite{izacard2021fid}, while Self-RAG integrates retrieval and self-reflection decisions \cite{asai2024selfrag}. These architectures motivate evidence-aware generation, but a displayed citation may still be irrelevant to the nearby claim, and a source-supported sentence may fail to answer the question.

RAGAS and ARES propose automated dimensions and synthetic workflows for RAG evaluation \cite{es2024ragas,saadfalcon2024ares}. FActScore decomposes long-form text into atomic factual claims \cite{min2023factscore}. SummEval and G-Eval illustrate both the need for multidimensional review and the limitations of automatic or model-based judges \cite{fabbri2021summeval,liu2023geval}. Following this literature, the thesis separates four constructs: whether a claim is supported somewhere in returned evidence; whether its attached citation is the correct evidence; whether expected answer points are covered; and whether the system declines questions outside the available evidence. Citation presence is recorded but is not allowed to stand in for the other three constructs.

Attribution also has an interface dimension. A paper ID without a passage requires the user to search the source again; a snippet without page or section context can be difficult to verify; and a ``grounded'' badge can imply more than a structural check established. The platform therefore returns paper, chunk, page, section, snippet, provider, model, warning, and review state. The later atomic review is intentionally stricter than the runtime structural verifier. This difference is not an inconsistency: the runtime can check locator integrity deterministically, while entailment and usefulness require a separate evidence process.

\section{Abstention and selective response}

Selective classification formalizes a reject option as a risk--coverage trade-off: a system may reduce error on answered cases by declining cases for which evidence is insufficient \cite{elyaniv2010selective}. A RAG answerer is not identical to a classifier, so the theory is used here as a design analogy rather than a direct guarantee. The relevant principle is that answering every request is not automatically desirable when the corpus is bounded.

Abstention in research QA can fail at several levels. A query may be outside the catalogue; relevant papers may exist but lack eligible text; retrieval scores may be weak or contradictory; retrieved passages may not cover the requested answer points; or a provider may be unavailable. These conditions should not collapse into one generic error. The interface can return an explicit out-of-corpus notice, evidence-only results, a provider-fallback warning, or no generated answer. Evaluation must then report both coverage and error among answered cases, with enough unanswerable cases to estimate the trade-off. The present small unanswerable subsets are diagnostic, not sufficient for calibration.

\section{Recommendation and human-centered advice}

Research-paper recommendation combines content, citations, profiles, and hybrid signals; surveys identify reproducible evaluation and realistic user modeling as continuing challenges \cite{beel2016researchpaper}. Publication history can represent recent research interests \cite{sugiyama2010scholarly}. Broader recommender evaluation literature emphasizes that offline predictive measures, user tasks, and whole-system/user evaluation answer different questions \cite{herlocker2004evaluating}. A ranked paper may be relevant without producing a feasible project, and a well-formed idea may still be unsuitable for a particular programme or supervisor.

This distinction is central to the Thesis Extension Finder. Its profile contains interests, skills, available time, project type, data constraints, difficulty, preferred topics, and avoid-topics. Those inputs can support candidate retrieval and transparent heuristics. They cannot verify data permission, novelty, prerequisite mastery, institutional fit, or supervision capacity. The design therefore separates paper-supported baseline work; explicit paper-stated future work; an inferred or missing gap; and a system-generated suggestion. It also offers an evidence-only mode. The full output structures an MVP, stretch goals, risks, skills, data assumptions, and evaluation plan, but each remains advice requiring external confirmation.

Human-centered AI guidance recommends making capabilities and limits clear, supporting correction, showing contextually relevant information, and enabling user control \cite{amershi2019guidelines}. Interactive machine-learning work likewise emphasizes user feedback throughout the system lifecycle \cite{amershi2014power}. Applied here, those principles favor profile echoing, editable inputs, visible assumptions, evidence links, alternative evidence-only output, error recovery, and reviewer correction over a single opaque suitability score. The present study tests those contracts as software behavior; it does not substitute automated accessibility checks or AI-proxy ratings for a human usability study.

\section{Topics, authors, and identity boundaries}

Latent Dirichlet allocation represents documents as distributions over latent topics \cite{blei2003lda}. BERTopic combines transformer embeddings, clustering, and class-based term weighting \cite{grootendorst2022bertopic}. Topic interpretability remains difficult to reduce to one numerical measure \cite{doogan2021topic}. For a small curated corpus, the project retains an editable controlled vocabulary and evaluates it against a dense label-prototype alternative. The controlled method trades open-ended discovery for determinism and governance; precision, recall, coverage, and per-label support must determine how much that trade achieved.

Expert finding has been formalized through candidate- and document-centric language models \cite{balog2006formal,balog2009language}. In a public laboratory setting, however, an author--topic score may be misread as current expertise, availability, or willingness to supervise. The platform restricts author-topic evidence to eligible indexed publications, preserves aliases and unresolved candidates, and states that the relation is publication-derived. Identity merges require external confirmation because a false merge changes attribution across every dependent view.

\section{Review, documentation, and responsible AI}

Datasheets and model cards formalize provenance, intended use, evaluation, and limitation reporting \cite{gebru2021datasheets,mitchell2019modelcards}. Language-model risk taxonomies include misinformation and interaction harms \cite{weidinger2022taxonomy}; the NIST AI Risk Management Framework organizes governance around mapping, measurement, management, and oversight \cite{tabassi2023airmf}. These frameworks do not prescribe this application's exact schema, but they motivate persistent source/configuration metadata, declared intended use, failure reporting, and accountable review.

The platform operationalizes those ideas through generated-content notices, actor-typed review states, authenticated mutations, append-only event hashes, transient public profiles, and a sanitized release. A correction is not a silent overwrite: the event identifies the target, prior and new state, actor role/type, rationale, time, and chain link. An AI reviewer can create an AI-reviewed state but cannot issue a human approval state. Content that cannot be reproduced safely remains marked for reprocessing rather than being deleted or relabelled as current.

\section{Synthesis and research gap}

\Cref{tab:literature-synthesis} connects the literature-derived risk to the implemented response and the evidence still required. The table is a positioning device, not a capability-superiority claim.

\begin{table}[htbp]
\centering
\small
\caption{Literature-derived design tensions and this thesis's bounded response.}
\label{tab:literature-synthesis}
\begin{tabularx}{\textwidth}{P{35mm}P{48mm}Y}
\toprule
Design tension & Implemented response & Evidence boundary \\
\midrule
Corpus scale versus inspectable provenance & page-aware local corpus; eligibility states; hash-bound indexes & internal eligible-set completeness, not global scholarly coverage \\
Lexical versus dense/hybrid retrieval & common candidate/evaluation boundary; development tuning; ablation & small source-derived AI-silver test, not universal ranking superiority \\
Citation presence versus attribution & atomic support and claim--citation review; answer points & AI-reviewed silver evidence; no independent human reliability \\
Answer coverage versus abstention & distinct answerability and unanswerable outcomes; visible fallbacks & diagnostic unanswerable set, not calibrated selective risk \\
Paper relevance versus project suitability & evidence-only baseline; fact/gap/suggestion separation; explicit assumptions & synthetic profiles and AI-proxy review, not student/supervisor benefit \\
Topic/author discovery versus over-attribution & controlled vocabulary; per-label evidence; conservative identities & publication-derived association, not availability or endorsement \\
Generated output versus accountable correction & actor-typed states and append-only review events & implemented auditability, not completed institutional governance \\
\bottomrule
\end{tabularx}
\end{table}

Prior work supplies mature components for scholarly graphs, retrieval, recommendation, RAG, topic models, human--AI interaction, and documentation. The gap addressed here is their evidence-calibrated integration for a small laboratory: full-paper ingestion; paper/chunk/page/section traceability; source-fact, future-work, inference, and suggestion separation; evidence-only and generative recommendation modes; attributable review; and failure-oriented evaluation. The contribution is an engineered and evaluated integration, not a new foundational model, a complete bibliographic index, or a validated educational intervention.
